\section{A uniform local law}\label{sec:llt}

So far we compared each bin with an endpoint. Now we approximate $P_N(k)$
itself by a single formula. This is a local form of Kagey's normal limit, but
we want it uniform in $q\le\frac12$: in the diffusive regime ($q$ fixed) the
bin is close to Gaussian, while in the ballistic regime ($T=Nq$ bounded) the
walk switches only a bounded number of times and the Gaussian form fails
(Proposition~\ref{prop:llt-gauss}). We first write $P_N(k)$ exactly through a
difference of 2 binomial variables, then replace that difference by a Skellam
law with the same mean and variance. Throughout, $n_1=k-1$, $n_2=N-k-1$ and
$n_{\min}=\min(n_1,n_2)$.

\begin{theorem}[a difference of 2 binomials]\label{thm:llt-identity}
Let $N\ge2$, $1\le k\le N-1$ and $0<q<1$. Let $Y_1\sim\mathrm{Bin}(n_1,q)$ and
$Y_2\sim\mathrm{Bin}(n_2,q)$ be independent, and put $D=Y_1-Y_2$. Then
\begin{equation}\label{eq:llt-identity}
P_N(k)=\frac q2\bigl[2\Prob(D=0)+\Prob(D=1)+\Prob(D=-1)\bigr].
\end{equation}
\end{theorem}

In words, $Y_1+1$ and $Y_2+1$ play the roles of the numbers of runs of
$\sR$'s and of $\sL$'s. In a word these differ by at most 1, which is why only
$D\in\{-1,0,1\}$ appears, and the factor $q$ pays for the 1 switch that every
interior bin needs.

\begin{proof}
We rewrite Theorem~\ref{thm:pmf}, where $a-1=n_1$ and $b-1=n_2$. Note that
$\Prob(Y_1=r)\Prob(Y_2=r')=\binom{n_1}r\binom{n_2}{r'}q^{r+r'}p^{N-2-r-r'}$.
The term $j=2r+1$ of Theorem~\ref{thm:pmf} is
$\binom{n_1}r\binom{n_2}rp^{N-2-2r}q^{2r+1}=q\Prob(Y_1=r)\Prob(Y_2=r)$, and
summing over $r$ gives $q\Prob(D=0)$. The term $j=2r+2$ is
\[
\tfrac12\Bigl[\tbinom{n_1}{r+1}\tbinom{n_2}r+\tbinom{n_1}r\tbinom{n_2}{r+1}\Bigr]p^{N-3-2r}q^{2r+2}
=\tfrac q2\bigl[\Prob(Y_1=r+1)\Prob(Y_2=r)+\Prob(Y_1=r)\Prob(Y_2=r+1)\bigr],
\]
and summing over $r$ gives $\frac q2[\Prob(D=1)+\Prob(D=-1)]$.
\end{proof}

For $k=1$ we have $n_1=0$ and $D=-Y_2$, so $\Prob(D=0)=p^{N-2}$,
$\Prob(D=-1)=(N-2)qp^{N-3}$ and $\Prob(D=1)=0$. By \eqref{eq:llt-identity} and
symmetry,
\begin{equation}\label{eq:llt-edge}
P_N(1)=P_N(N-1)=qp^{N-2}+\tfrac12(N-2)q^2p^{N-3}.
\end{equation}
With $P_N(0)=P_N(N)=\frac12p^{N-1}$, the 4 bins $0,1,N-1,N$ are exact, and from
now on $N\ge4$ and $2\le k\le N-2$, so that $n_1,n_2\ge1$.

Next we tilt $D$ to mean $0$. For $s>0$ put
\[
r_1=\frac{qs}{p+qs},\qquad r_2=\frac q{ps+q},\qquad v_i=r_i(1-r_i),\qquad
\Lambda(s)=(p+qs)^{n_1}(p+q/s)^{n_2}.
\]
Under the tilted law $\Prob_s$, $Y_1\sim\mathrm{Bin}(n_1,r_1)$ and
$Y_2\sim\mathrm{Bin}(n_2,r_2)$ are independent, and
$\Prob(D=j)=\Lambda(s)s^{-j}\Prob_s(D=j)$ (Lemma~\ref{lem:llt-tilt}). There is
a unique $s>0$ with $n_1r_1=n_2r_2=:m$, so that $\E_sD=0$
(Lemma~\ref{lem:llt-saddle}). We fix this $s$, let $x=n_1v_1+n_2v_2$ be the
variance of $D$ under $\Prob_s$, and define
\begin{equation}\label{eq:llt-U}
U_N(k)=\frac q2\Lambda(s)e^{-x}\bigl[2I_0(x)+(s+s^{-1})I_1(x)\bigr].
\end{equation}
At the center of an even row, $n_1=n_2$ gives $s=1$, $r_1=r_2=q$ and
$\Lambda=1$, so
\begin{equation}\label{eq:llt-center}
U_N(N/2)=q\,e^{-x}\bigl(I_0(x)+I_1(x)\bigr),\qquad x=(N-2)pq.
\end{equation}

Here is what \eqref{eq:llt-U} does. By \eqref{eq:llt-identity} and the tilt,
$P_N(k)=\frac q2\Lambda(s)\bigl[2\Prob_s(D=0)+s^{-1}\Prob_s(D=1)+s\Prob_s(D=-1)\bigr]$.
The law $e^{-x}I_{|j|}(x)$ on $\Z$ is the symmetric Skellam law, the law of a
difference of 2 independent Poisson variables with mean $x/2$, and it has mean
$0$ and variance $x$. So $U_N(k)$ replaces the tilted $D$ by the Skellam law
with the same mean and variance. This law is exact for a Poisson difference,
which $D$ resembles when $q$ is small, and close to Gaussian when $x$ is large.
At the center $D$ is a sum of $N/2-1$ independent variables with
$\Prob(\pm1)=pq$, and replacing it by this Skellam law is the approximation
that \v{C}ekanavi\v{c}ius and Liaudanskait\.{e}~\cite[Eq.~(3)]{cekanavicius2024}
attribute to Presman. The tilt extends it off the center.

To state the error bound, let $\kappa=1-\cos t$ and write $\langle h\rangle$
for the average of $h(t)$ against the density proportional to $e^{x\cos t}$ on
$[-\pi,\pi]$. With $\mathcal R=\mathcal R(x)=I_1(x)/I_0(x)$,
Lemma~\ref{lem:llt-moments} gives
\[
\langle\kappa^2\rangle=2-2\mathcal R-\frac{\mathcal R}x,\qquad
\langle\kappa^3\rangle=4-4\mathcal R-\frac{3\mathcal R}x-\frac{2\mathcal R}{x^2}+\frac1x.
\]
Put $V=n_1v_1^2+n_2v_2^2$ and $M=m\max(r_1,r_2)$.

\begin{theorem}[uniform local law]\label{thm:llt}
Let $N\ge4$, $2\le k\le N-2$ and $0<q<1$. Then $P_N(k)=U_N(k)(1+\theta)$ with
\[
|\theta|\le\mathcal E:=\frac1{\mathcal R(x)}\Bigl[\Bigl(4V+\frac{64}3M\Bigr)\langle\kappa^2\rangle
+\frac{1024}9M^2\langle\kappa^3\rangle\Bigr].
\]
\end{theorem}

In words, the relative error is controlled by 2 things: how far the tilted
binomial steps are from Poisson ($V$, which is small when the $r_i$ are small)
and how skewed they are ($M$), each weighted by a moment of $\kappa$ that is
small when the variance $x$ is large.

\emph{Idea of the proof.} By Fourier inversion, $P_N(k)$ and $U_N(k)$ are
$\frac q2\Lambda(s)\frac1{2\pi}\int_{-\pi}^{\pi}f\omega\,dt$ with
$\omega(t)=2+se^{it}+s^{-1}e^{-it}$, where $f$ is $\chi(t)=\E_se^{itD}$ for
$P_N(k)$ and the Skellam characteristic function $g(t)=e^{x(\cos t-1)}$ for
$U_N(k)$. Since $|1-r+re^{it}|^2=1-2v\kappa=e^{-2v\kappa}(1+O(v^2\kappa^2))$
with $v=r(1-r)$, we get $|\chi|\le g$ and $g-|\chi|=O(V\kappa^2)\,g$. The phase
of $\chi$ has no linear term since $\E_sD=0$, so it is $O(M|t|^3)$. Hence
$|\theta|$ is at most a constant times
$\bigl((V+M)\langle\kappa^2\rangle+M^2\langle\kappa^3\rangle\bigr)/\mathcal R(x)$,
where $\langle\kappa^j\rangle=O(\min(1,x^{-j}))$ and
$1/\mathcal R(x)\le1+2/x$. Appendix~\ref{app:llt} gives the proof with
explicit constants.

\begin{corollary}\label{cor:llt}
Let $N\ge4$, $2\le k\le N-2$ and $0<q\le\frac12$, and put
$r_{\max}=\max(r_1,r_2)$. Then
\[
|\theta|\le\min\Bigl\{\frac{22334}{n_{\min}},\ 221\,r_{\max}+5825\,r_{\max}^2\Bigr\}.
\]
Also $r_{\max}\le\delta_T:=2T/\sqrt N+4T^2/N$, so
$|\theta|\le221\,\delta_T+5825\,\delta_T^2$ for every bin $2\le k\le N-2$. At
the center of an even row, $r_{\max}=q$ and $n_{\min}=N/2-1$.
\end{corollary}

So the relative error tends to $0$ uniformly in $0<q\le\frac12$ as
$n_{\min}\to\infty$, and uniformly over $2\le k\le N-2$ as $T/\sqrt N\to0$,
which covers the ballistic regime $T=O(1)$ up to the edges. The constants are
crude. The largest $n_{\min}|\theta|$ we observed is $0.686$ (at $N=10^4$,
$q=\frac12$, over $2\le k\le400$ with $q=\frac12$ or $0.3$, and over every bin
of our grid with $N\le1000$ and $10^{-5}\le q\le\frac12$), while the first
bound has the constant $22334$ and exceeds $1$ unless $n_{\min}>22334$. Near an edge the largest error we
observed is $|\theta|=0.1617$ ($k=3$, $N=10^4$, $q=\frac12$), an observation
and not a bound. For $q>\frac12$ Theorem~\ref{thm:llt} still holds, but
$\mathcal E$ need not be small, and a uniform form is open. Near $q=1$ the
symmetric kernel fails once $n_1\ne n_2$. We found $\theta=0.578$ at the
middle bin $k=50$ of the odd row $N=101$ with $q=0.999$, while $\theta\approx0$
at $N=100$, $k=50$.

Why not a Gaussian form? Replacing $e^{-x}I_j(x)$ in \eqref{eq:llt-U} by
$(2\pi x)^{-1/2}$ gives the Gaussian saddle point
$G_N(k)=\frac q2\Lambda(s)(2+s+s^{-1})(2\pi x)^{-1/2}$. We also define the
transfer-matrix form $\mathrm{TM}_N=2(2\pi(N-1)p/q)^{-1/2}$, the Gaussian local
form for $\Prob(X_N=0)$ with the variance $(N-1)p/q$ that $\lambda_+^{N-1}$
gives in the proof of Theorem~\ref{thm:clt}. The factor $2$ is the span of the
lattice of $X_N=2k-N$.

\begin{proposition}[Gaussian saddle points]\label{prop:llt-gauss}
Fix an even $N\ge4$. As $q\to0$,
\[
\frac{P_N(N/2)}q\to1,\qquad
\frac{G_N(N/2)}{P_N(N/2)}\sim\frac2{\sqrt{2\pi(N-2)q}}\to\infty,\qquad
\frac{\mathrm{TM}_N}{P_N(N/2)}\sim\frac2{\sqrt{2\pi(N-1)q}}\to\infty.
\]
\end{proposition}

\begin{proof}
The term $j=1$ of Theorem~\ref{thm:pmf} is $qp^{N-2}$ and the other terms are
$O(q^2)$, so $P_N(N/2)=q+O(q^2)$. At the center $s=1$, $\Lambda=1$ and
$x=(N-2)pq$, so $G_N(N/2)=2q/\sqrt{2\pi(N-2)pq}$. Also
$\mathrm{TM}_N=2q/\sqrt{2\pi(N-1)pq}$. Let $q\to0$.
\end{proof}

More generally, since $P_N(N/2)/q\to1$, any approximation $A_N(q)$ of the
center with $A_N(q)/q\to\infty$ has relative error $P_N(N/2)/A_N(q)-1\to-1$.
Both saddle points are of this kind. We make no claim about other Gaussian
forms. The Bessel form avoids this because $e^{-x}I_j(x)$ is exact for the
Poisson difference that $D$ approaches in the ballistic regime.

\begin{remark}[the center crossing]\label{rem:llt-center}
For even $N\ge4$ let $\mathcal G(q)$ be the logarithm of the ratio of
$U_N(N/2)$ to $P_N(0)=\frac12p^{N-1}$. By $I_0'=I_1$, $I_1'=I_0-I_1/x$
\cite[Eqs.~10.29.2, 10.29.3]{dlmf} and $I_1\le I_0$, we get
$\mathcal G'(q)\ge\frac1{2pq}+\frac{N-1}p>N-1$ on $(0,\frac12]$. Also
$\mathcal G\to-\infty$ as $q\to0$ and
$\mathcal G(\frac12)\ge(N-1)\log2-\frac{N-2}4>0$, so $\mathcal G$ has a unique
root $q_U$ there, an approximate crossing. In Table~\ref{tab:llt-center},
which is purely numerical, the relative error of $q_U$ is 16 to 36 times
smaller than that of the Bessel root $u_B$ of Theorem~\ref{thm:bessel}. Put
$\vartheta_N=221q_N+5825q_N^2$. If $\vartheta_N<1$, Corollary~\ref{cor:llt}
gives $|\mathcal G(q_N)|\le\vartheta_N/(1-\vartheta_N)$, and the mean value
theorem gives $|q_U-q_N|\le\vartheta_N/((1-\vartheta_N)(N-1))$. By our
computation $\vartheta_N\ge1$ for every even $N<1454$, and beyond it this bound
is weaker than Theorem~\ref{thm:bessel}(ii) at every $N$ we computed.
\end{remark}

\begin{table}[ht]
\centering\small
\begin{tabular}{@{}rcccc@{}}
\toprule
$N$ & $p_N$ & $q_U/q_N-1$ & $u_B/u_N-1$ & ratio\\
\midrule
100 & 0.9649663263 & $-1.477\times10^{-3}$ & $-2.39\times10^{-2}$ & 0.062\\
400 & 0.9881246429 & $-3.083\times10^{-4}$ & $-7.61\times10^{-3}$ & 0.041\\
1600 & 0.9962346137 & $-6.469\times10^{-5}$ & $-2.31\times10^{-3}$ & 0.028\\
\bottomrule
\end{tabular}
\caption{The root $q_U$ of $U_N(N/2)=P_N(0)$ against the Bessel root $u_B$ of
Theorem~\ref{thm:bessel}. The ratio is $(q_U/q_N-1)/(u_B/u_N-1)$. Numerical values, not bounds.}
\label{tab:llt-center}
\end{table}

\paragraph{Related local limits.} \v{C}ekanavi\v{c}ius and
Liaudanskait\.{e}~\cite{cekanavicius2024} and \v{C}ekanavi\v{c}ius and
Vellaisamy~\cite{cekanavicius2010} give total-variation, local and Wasserstein
bounds with unspecified absolute constants for chains outside our regime. One
is a symmetric 3-state chain whose step is $0$ with probability at least
$\frac{14}{15}$. The other is a 2-state chain with $\Prob(1\mid1)\le\frac12$,
and some of its bounds also need $\Prob(1\mid0)\le\frac1{30}$. Neither allows
both states to persist, as our $p\to1$ does. Jensen~\cite{jensen2013}
conditions the count $S_n$ on $0<S_n<n$ and patches small tilted variances,
and his relative error stays bounded (about $0.18$, found partly by numerics)
but does not tend to $0$. In the literature we searched we did not find a
tilted Skellam kernel off the center with a relative local error, explicit
constants and uniformity in $q\le\frac12$.
